% ==========================================================================
%  Section 13 -- What is machine-checked and what is not.  Owner: W10.
%  Frozen label: sec:lean-boundary.
%
%  Architecture: a numbered list of divergences, each a fact about the sources,
%  with no qualifier that softens it.  Every item is backed by notes/lean_ledger.md
%  sections 4 and 5; the counts that exist as macros are macros, the reachability
%  count is the word "eight" (CONTRACT section 11, final paragraph), and the direct
%  import count is deliberately not stated (it would be a typed structural count).
%
%  Traps observed here.  (1) The root aggregator does not build every module, and
%  the footprints of six registered theorems rest on the transcript, not a build.
%  (2) The final-theorem wrapper proves nothing; it may be named, and must be
%  called vacuous in the same breath.  (3) The total-lift zero convention is not a
%  null-set assertion; the chain avoids it.  (4) The measure chain is 2-dimensional
%  while the covariance theorems are general-d.
% ==========================================================================
\section{What is machine-checked and what is not}\label{sec:lean-boundary}

The formal side of this paper is easier to overstate than to state, so the
divergences between the sources and a maximally favourable reading are collected
here as numbered facts. Each is a property of the sources, not a hedge on a
result: the declarations asserted below are the declarations that exist, and the
sentences below say how far the evidence for them extends. The influence of each
item on the mathematics is none; the influence on what a reader may conclude is
total.

\begin{enumerate}

\item \emph{Reachability.} \NumCleanModules{} modules exist in the package, and
  the package root aggregator imports a subset of them directly.
  \NumUnreachableModules{} modules are not reachable from the root file even
  transitively. Two of them are load-bearing:
  \lmod{PeriodicHMinusThreeWeakLimitMarginal}, which declares
  \NumUnreachableChainTheorems{} of the \NumMeasureContractTheorems{} registered
  measure-limit theorems, and
  \lmod{PeriodicHMinusThreeProkhorovSubsequence}, which declares the Prokhorov
  extraction step of the argument. Consequently a build of the \emph{root file}
  does not evidence the axiom footprints of those declarations: nothing reachable
  from it contains them, so a claim that the root covers the development is a
  claim about a subset of the package. What does evidence them is a build of the
  library target, which elaborates every module including the unreachable ones,
  and that is the build Appendix~B records. The distinction matters to a reader
  who builds the aggregator and stops: the result is green and says nothing about
  \NumUnreachableModules{} modules, \NumUnreachableChainTheorems{} of whose
  declarations are on this paper's chain.

\item \emph{A vacuous wrapper.} The development contains a
  final-theorem-shaped interface,
  \lmod{latticeFreeFieldMeasure_tendsto_periodicFreeGaussian}, whose proof is
  \texttt{exact hfinal}: it assumes its own conclusion. This paper does not use
  it, and no result on these pages depends on it. It exists as a contract-shaped
  interface, listing the data that a completed assembly would have to supply --
  tightness, the finite marginals, the cylinder-to-Borel input and the
  subsequential identification -- and a reader who found it by its name and its
  shape would take it for the paper's theorem. It is not, and citing it as such
  would be filing an empty citation.

\item \emph{The formalisation's lift is not the paper's}. The development's
  carrier-valued lift is a total function, defined coordinatewise with the
  series appearing only in the $\ell^2$ side condition, together with the
  convention that it is $0$ off the convergence set; its measurability is one
  conjunct of an input of the formal development, the other conjunct asserting
  almost-everywhere agreement with the raw weighted lift. That convention is
  \emph{weaker} than a range statement and must not be read as one: the function
  is total at the price of the convention, not of a theorem about it. Nothing in
  the mathematics turns on it. The paper's carrier-valued representative is $L_N$
  of \cref{def:finite-source-lift}, whose total range, Borel measurability and
  entry formula are proved outright (\cref{lem:lift-representative}), and the
  chain of \cref{sec:main} runs through $L_N$, not through the development's
  lift; on the finite-source route the development moreover proves the two
  pushforwards equal without using the input. So a
  reader comparing the two sources should read the development's input as a fact
  about one implementation choice, not as a gap in \cref{thm:main}. What is
  genuinely absent on both sides is stated in \cref{sec:scope}.

\item \emph{Dimension.} The measure chain is hard-wired to two dimensions: the
  carrier, the Fourier index set, the cylinder maps and the law sequence are all
  built over the two-dimensional index set, and no parameter for the dimension
  occurs in them. The registered covariance theorems are stated for a general
  dimension $d$. The two statements are therefore not the same statement
  generalised differently: the covariance statements are \emph{registered} for
  general $d$ (that is how the Lean declarations quantify), this paper writes
  them in $d=2$, and the measure limit, as formalized, is available in $d=2$
  only. No dimension other than two is claimed anywhere in Part~I
  (\cref{rem:dimension}).

\item \emph{What a green build proves.} A complete build proves the Lean
  declarations that the sources state, under the hypotheses that the sources
  attach to them, on the toolchain recorded in the transcript; a complete build
  of this package exists and is recorded in \cref{sec:lean-reproduction}. It
  proves nothing
  about the mathematics of a hypothesis: a hypothesis is an input, and a build
  that consumes it is not evidence for it. It proves nothing about statements
  outside the library target of Item~1. And the count of \NumCleanSorry{}
  occurrences of \texttt{sorry} is a statement about the text of the sources,
  obtained by a scan that strips comments, rather than a statement about the
  mathematics: a source file free of placeholders can still formalize a
  proposition weaker than the one a reader has in mind, which is why Item~3 and
  Item~4 are stated here rather than absorbed into the count. The mathematical
  boundary of this paper -- the coverage comparison between these ten
  declarations and the rest of the development, and the absence of a rate -- is a
  different boundary, and \cref{sec:scope} draws it.

\end{enumerate}
